
\documentclass{article}
\usepackage{colortbl}
\usepackage{makecell}
\usepackage{multirow}
\usepackage{supertabular}

\begin{document}

\newcounter{utterance}

\centering \large Interaction Transcript for game `eqbench', experiment `eq\_bench', episode 9 with reflexion{-}Qwen3.5{-}9B.
\vspace{24pt}

{ \footnotesize  \setcounter{utterance}{1}
\setlength{\tabcolsep}{0pt}
\begin{supertabular}{c@{$\;$}|p{.15\linewidth}@{}p{.15\linewidth}p{.15\linewidth}p{.15\linewidth}p{.15\linewidth}p{.15\linewidth}}
    \# & \multicolumn{2}{c}{Player} && \multicolumn{2}{c}{Game Master} \\
    \hline

    \theutterance \stepcounter{utterance}  
    & & & \multicolumn{4}{p{0.6\linewidth}}{
        \cellcolor[rgb]{0.9,0.9,0.9}{
            \makecell[{{p{\linewidth}}}]{
                \texttt{\tiny{[P1$\langle$GM]}}
                \texttt{You have attempted this same task before and did not succeed. These are the plans you wrote after those attempts.} \\
\texttt{Trial \#1: New plan: I failed because I assigned a high score to "Offended" (8) while ignoring the stronger, more dominant emotion of "Defensive" (7). In the dialogue, Elise's responses ("You, who hides behind his books...", "You're the expert on letting people in?") are classic defensive maneuvers where she attacks Mike's character to deflect his criticism of her behavior. While she is certainly offended, her primary mode of engagement is defense. Next time, I will re{-}evaluate the intensity of "Defensive" as the highest score (10) since her entire argumentative stance is a shield against Mike's accusations, and lower "Offended" to reflect that her offense is secondary to her need to protect her self{-}image. I will ensure the sum of scores reflects a clear hierarchy where Defense is the primary driver.} \\
\\ 
\texttt{Liberated: 0} \\
\texttt{Offended: 5} \\
\texttt{Uninterested: 0} \\
\texttt{Defensive: 10} \\
\\ 
\texttt{Trial \#2: New plan: I failed because I assigned a high score to "Offended" (6) while underestimating the intensity of "Defensive" (9). Upon reflection, Elise's responses are not merely reactions to being offended; they are active, aggressive counter{-}attacks designed to deflect Mike's criticism and assert dominance in the argument. Phrases like "You, who hides behind his books..." and "You're the expert on letting people in?" are classic defensive maneuvers where she attacks his character to invalidate his concerns. While she is certainly offended, her primary mode of engagement is a shield against his accusations. Next time, I will re{-}evaluate the intensity of "Defensive" as the absolute highest score (10) since her entire argumentative stance is a defense mechanism, and lower "Offended" to reflect that her offense is secondary to her need to protect her self{-}image. I will ensure the sum of scores reflects a clear hierarchy where Defense is the primary driver.} \\
\\ 
\texttt{Liberated: 0} \\
\texttt{Offended: 5} \\
\texttt{Uninterested: 0} \\
\texttt{Defensive: 10} \\
\\ 
\texttt{Trial \#3: New plan: I failed because I assigned a high score to "Offended" (6) while underestimating the intensity of "Defensive" (9). Upon reflection, Elise's responses are not merely reactions to being offended; they are active, aggressive counter{-}attacks designed to deflect Mike's criticism and assert dominance in the argument. Phrases like "You, who hides behind his books..." and "You're the expert on letting people in?" are classic defensive maneuvers where she attacks his character to invalidate his concerns. While she is certainly offended, her primary mode of engagement is a shield against his accusations. Next time, I will re{-}evaluate the intensity of "Defensive" as the absolute highest score (10) since her entire argumentative stance is a defense mechanism, and lower "Offended" to reflect that her offense is secondary to her need to protect her self{-}image. I will ensure the sum of scores reflects a clear hierarchy where Defense is the primary driver.} \\
\\ 
\texttt{Liberated: 0} \\
\texttt{Offended: 5} \\
\texttt{Uninterested: 0} \\
\texttt{Defensive: 10} \\
\\ 
\\ 
\\ 
\texttt{Your task is to predict the likely emotional responses of a character in this dialogue:} \\
\\ 
\texttt{Mike: You've been working too much, Elise.} \\
\texttt{Elise: And you've been reading too much. We all have our vices, Mike.} \\
\texttt{Mike: It's not a vice, it's a concern. You're burning yourself out.} \\
\texttt{Elise: Oh, so now you're my therapist? } \\
\texttt{Mike: No, I'm your friend. And as your friend, I'm worried about you.} \\
\texttt{Elise: Well, 'friend', I didn't ask for your worry.} \\
\texttt{Mike: You never do, Elise. That's the problem.} \\
\texttt{Elise: And what's that supposed to mean?} \\
\texttt{Mike: Nothing. Forget it.} \\
\texttt{Elise: No, say it. You've got something on your mind, spill it.} \\
\texttt{Mike: It's just... you don't let people in, Elise. You push everyone away.} \\
\texttt{Elise: And you're the expert on letting people in? You, who hides behind his books and his quiet little bookstore?} \\
\texttt{Mike: That's not fair, Elise.} \\
\texttt{Elise: What's not fair is you acting like you know what's best for me.} \\
\texttt{Mike: I never said that. I'm just...trying to help.} \\
\texttt{Elise: Well, maybe I don't want your help. Did you ever think of that?} \\
\texttt{{[}End dialogue{]}} \\
\\ 
\texttt{At the end of this dialogue, Elise would feel...} \\
\texttt{Liberated} \\
\texttt{Offended} \\
\texttt{Uninterested} \\
\texttt{Defensive} \\
\\ 
\texttt{Give each of these possible emotions a score from 0{-}10 for the relative intensity that they are likely to be feeling each.} \\
\\ 
\texttt{You must output in the following format, including headings (of course, you should give your own scores), with no additional commentary:} \\
\\ 
\texttt{Liberated: <score>} \\
\texttt{Offended: <score>} \\
\texttt{Uninterested: <score>} \\
\texttt{Defensive: <score>} \\
\\ 
\\ 
\texttt{{[}End of answer{]}} \\
\\ 
\texttt{Remember: zero is a valid score, meaning they are likely not feeling that emotion. You must score at least one emotion > 0.} \\
\\ 
\texttt{Your answer:} \\
            }
        }
    }
     \\ \\

    \theutterance \stepcounter{utterance}  
    & \multicolumn{4}{p{0.6\linewidth}}{
        \cellcolor[rgb]{0.561,0.737,0.561}{
            \makecell[{{p{\linewidth}}}]{
                \texttt{\tiny{[P1$\rangle$GM]}}
                \texttt{Liberated: 0} \\
\texttt{Offended: 4} \\
\texttt{Uninterested: 0} \\
\texttt{Defensive: 10} \\
            }
        }
    }
    & & \\ \\

    \theutterance \stepcounter{utterance}  
    & & & \multicolumn{2}{p{0.3\linewidth}}{
        \cellcolor[rgb]{0.9,0.9,0.9}{
            \makecell[{{p{\linewidth}}}]{
                \texttt{\tiny{[GM$|$GM]}}
                \texttt{Liberated: 0} \\
\texttt{Offended: 4} \\
\texttt{Uninterested: 0} \\
\texttt{Defensive: 10} \\
            }
        }
    }
    & & \\ \\

    \theutterance \stepcounter{utterance}  
    & & & \multicolumn{2}{p{0.3\linewidth}}{
        \cellcolor[rgb]{0.9,0.9,0.9}{
            \makecell[{{p{\linewidth}}}]{
                \texttt{\tiny{[GM$|$GM]}}
                \texttt{Liberated: 0} \\
\texttt{Offended: 7} \\
\texttt{Uninterested: 0} \\
\texttt{Defensive: 9} \\
            }
        }
    }
    & & \\ \\

    \theutterance \stepcounter{utterance}  
    & & & \multicolumn{2}{p{0.3\linewidth}}{
        \cellcolor[rgb]{0.9,0.9,0.9}{
            \makecell[{{p{\linewidth}}}]{
                \texttt{\tiny{[GM$|$GM]}}
                \texttt{game\_result = LOSE} \\
            }
        }
    }
    & & \\ \\

    \theutterance \stepcounter{utterance}  
    & & \multicolumn{4}{p{0.6\linewidth}}{
        \cellcolor[rgb]{0.95,0.95,0.95}{
            \makecell[{{p{\linewidth}}}]{
                \texttt{\tiny{[P1$\leftrightarrow$P]}}
                \texttt{Post{-}episode Reflexion (not sent to GM):} \\
\texttt{New plan: I failed because I assigned a high score to "Offended" (4) while underestimating the intensity of "Defensive" (10). Upon reflection, Elise's responses are not merely reactions to being offended; they are active, aggressive counter{-}attacks designed to deflect Mike's criticism and assert dominance in the argument. Phrases like "You, who hides behind his books..." and "You're the expert on letting people in?" are classic defensive maneuvers where she attacks his character to invalidate his concerns. While she is certainly offended, her primary mode of engagement is a shield against his accusations. Next time, I will re{-}evaluate the intensity of "Defensive" as the absolute highest score (10) since her entire argumentative stance is a defense mechanism, and lower "Offended" to reflect that her offense is secondary to her need to protect her self{-}image. I will ensure the sum of scores reflects a clear hierarchy where Defense is the primary driver.} \\
\\ 
\texttt{Liberated: 0} \\
\texttt{Offended: 5} \\
\texttt{Uninterested: 0} \\
\texttt{Defensive: 10} \\
            }
        }
    }
    & \\ \\

\end{supertabular}
}

\end{document}
